\documentclass[11pt]{article}
\usepackage[a4paper, top=18mm, bottom=18mm, left=20mm, right=20mm]{geometry}
\usepackage[T2A]{fontenc}
\usepackage[utf8]{inputenc}
\usepackage[ukrainian]{babel}
\usepackage{graphicx}
\setlength{\parindent}{0pt}
\setlength{\parskip}{4pt}
\pagestyle{empty}
\begin{document}
%begin
{\large\textbf{Контрольна робота}}

\textbf{Варіант \No\,2}\hfill \textit{ПІБ:}\,\underline{\hspace{60mm}}
\vspace{4mm}

%beginitem
1. Що буде виведено за виконання фрагменту коду?
%code

try:
    print(3, end="")
    print(int(6//0.0), end="")
    print(7, end="")
except TypeError: 
    print(1, end="")
except Exception: 
    print(8, end="")
else:
    print(4, end="")
finally:
    print(0, end="")

%code

%enditem

%beginitem
2. Що буде виведено за виконання фрагменту коду?
try:
    res = 3**False-2
    if isinstance(res, bool):
        print(f'bool;{res}')
    elif isinstance(res, int):
        print(f'int;{res}')
    elif isinstance(res, float):
        print(f'float;{res:.2f}')
    else:
        print('???')
except:
    print('Z')

%enditem

%beginitem
3. Що буде виведено за виконання фрагменту коду?
a,b,c=3,6,7
    def g(b):
      global c      a=2
      b=1
      c=4
      return a+b+c
    a,b,c=1,8,4
    print(g(b),a,b,c)
%enditem

%beginitem
4. Що буде виведено за виконання фрагменту коду?
%code

try:
    print(3,end="")
    print(int(6//0.0),end="")
    print(7,end="")
except TypeError: print(1,end="")
except Exception: print(8,end="")
else: print(4,end="")
finally: print(0,end="")

%code

%enditem

%beginitem
5. %goodpath:Scripts/2019/Mod2019_1/Lex/03-good.txt
Відмітити все, що є коректним оператором С++:
( 1 ) n=2; n--

( 2 ) x=1//3

( 3 ) x in not y

( 4 ) x<y<z

%enditem

%beginitem
6. Що буде виведено за виконання фрагменту коду?
Записати ВИРАЗ, значенням якого є список, що складається з цілих чисел від 11 до 2016 (включно), причому спочатку за спаданням записано всі, що діляться на 7, а потім решту за спаданням.

%enditem

%beginitem
7. Що буде виведено за виконання фрагменту коду?
Рядки журналу через = містять ім'я користувача (ідентифікатор), час події,тип події, код події, наприклад
\begin{verbatim}
s="username = 2022-05-05 13:26 = ERROR = 404"
\end{verbatim}
у коді 
\begin{verbatim}
res = re.fullmatch('???', s) 
s = ''
code_str = ??? 
\end{verbatim}
чим треба замінити кожен з ???, щоб значенням code\_str був код події? 



%enditem

%beginitem
8. %goodpath:Scripts/2018/Mod2018_1/03-good.txt
Відмітити все, що є коректним оператором С++:
( 1 ) x=2'*'7;

( 2 ) int f(int n);

( 3 ) "a"=6;

( 4 ) x=y=z

%enditem

%beginitem
9. Що буде виведено за виконання фрагменту коду?
Записати ВИРАЗ, значенням якого є множина, що складається з цілих чисел від 11 до 2016 (включно), що діляться на 7.
%enditem

%beginitem
10. %goodpath:Scripts/2018/Mod2018_1/01-good.txt
Відмітити все, що є однією правильною лексемою:
( 1 ) <>

( 2 ) FIFO

( 3 ) 2*6

( 4 ) True

%enditem

%end
\newpage
\end{document}

